\begin{exercises}
  \item 该矩阵有两个特征值 $\lambda_1=3$、$\lambda_2=-4$。
    \begin{equation*}
      \begin{mat}
        4  &1 \\
        -8 &-5
      \end{mat}
    \end{equation*}
    求出两个不同的对角形矩阵，使它与它们相似。
    \begin{answer}
      我们可以调换列的位置，这只需调换做表示时所用的基。
      \begin{equation*}
        \begin{mat} 
          3 &0 \\ 
          0 &-4
        \end{mat}
        \qquad
        \begin{mat} 
          -4 &0 \\ 
          0 &3
        \end{mat}
      \end{equation*}
    \end{answer}
  \item 
    分别求出特征多项式与特征值。
    \begin{exparts*}
      \partsitem \( \begin{mat}[r]
                 10  &-9 \\
                  4  &-2
            \end{mat}  \)
      \partsitem $\begin{mat}[r]
                    1  &2  \\
                    4  &3  
                 \end{mat}$
      \partsitem \( \begin{mat}[r]
                  0  &3  \\
                  7 &0
                 \end{mat} \)
      \partsitem \( \begin{mat}[r]  
                  0  &0  \\
                  0  &0
            \end{mat}  \)
      \partsitem \( \begin{mat}[r]
                  1  &0  \\
                  0  &1
            \end{mat}  \)
    \end{exparts*}
    \begin{answer}
      \begin{exparts}
         \partsitem 这个
           \begin{equation*}
             0=
             \begin{vmatrix}
               10-x  &-9  \\
               4     &-2-x
             \end{vmatrix}
             =(10-x)(-2-x)-(-36)
           \end{equation*}
           化简后得到特征方程 \( x^2-8x+16=0 \)。 
           由于该方程可分解为 $(x-4)^2$，所以只有一个特征值
           \( \lambda_1=4 \)。
         \partsitem $0=(1-x)(3-x)-8=x^2-4x-5$；$\lambda_1=5$，$\lambda_2=-1$
         \partsitem \( x^2-21=0 \)； 
           \( \lambda_1=\sqrt{21} \)，$\lambda_2=-\sqrt{21}$
         \partsitem \( x^2=0 \)；\( \lambda_1=0 \)
         \partsitem \( x^2-2x+1=0 \)；\( \lambda_1=1 \)
       \end{exparts}  
     \end{answer}
  \recommended \item
    对每个矩阵，求出特征方程以及特征值和相应的特征向量。
    \begin{exparts*}
      \partsitem \( \begin{mat}[r]
                  3  &0  \\
                  8  &-1
            \end{mat}  \)
      \partsitem \( \begin{mat}[r]
                  3  &2  \\
                 -1  &0
            \end{mat}  \)
    \end{exparts*}
    \begin{answer}
       \begin{exparts}
         \partsitem 特征方程是 \( (3-x)(-1-x)=0 \)。
           它的根即特征值，为 \( \lambda_1=3 \) 与 
           \( \lambda_2=-1 \)。
           为求特征向量，我们考虑这个方程。
           \begin{equation*}
             \begin{mat}
               3-x  &0    \\
               8    &-1-x
             \end{mat}
             \colvec{b_1  \\  b_2}
             =\colvec[r]{0  \\  0}
           \end{equation*}
           为求与 $\lambda_1=3$ 相对应的特征向量，
           我们考虑由此得到的线性方程组。
           \begin{equation*}
             \begin{linsys}{2}
               0\cdot b_1  &+  &0\cdot b_2  &=  &0  \\
               8\cdot b_1  &+  &-4\cdot b_2 &=  &0
             \end{linsys}
           \end{equation*}
           特征子空间是第二个分量等于第一个分量两倍的向量组成的集合。
           \begin{equation*}
             \set{\colvec{b_2/2 \\ b_2}\suchthat b_2\in\C}
             \qquad
             \begin{mat}[r]
               3  &0  \\
               8  &-1
             \end{mat}
             \colvec{b_2/2  \\ b_2}
             =3\cdot\colvec{b_2/2 \\ b_2}
           \end{equation*}
           （这里取 $b_2$ 作为参数，只是因为它是上面方程组中的自由变量。）
           因此，这是一个与特征值 $3$ 相对应的特征向量。
           \begin{equation*}
             \colvec[r]{1 \\ 2}
           \end{equation*}

           求与 $\lambda_2=-1$ 相对应的特征向量与此类似。方程组
           \begin{equation*}
             \begin{linsys}{2}
               4\cdot b_1  &+  &0\cdot b_2  &=  &0  \\
               8\cdot b_1  &+  &0\cdot b_2 &=  &0
             \end{linsys}
           \end{equation*}
           导出由第一个分量为零的向量组成的集合。
           \begin{equation*}
             \set{\colvec{0 \\ b_2}\suchthat b_2\in\C}
             \qquad
             \begin{mat}[r]
               3  &0  \\
               8  &-1
             \end{mat}
             \colvec{0  \\ b_2}
             =-1\cdot\colvec{0 \\ b_2}
           \end{equation*}
           于是这是与 $\lambda_2$ 相对应的一个特征向量。
           \begin{equation*}
             \colvec[r]{0 \\ 1}
           \end{equation*}
         \partsitem 特征方程是
           \begin{equation*}
             0=
             \begin{vmatrix}
               3-x  &2  \\
               -1   &-x               
             \end{vmatrix}
             =x^2-3x+2=(x-2)(x-1)
           \end{equation*}
           因此特征值为 $\lambda_1=2$ 和 $\lambda_2=1$。
           为求特征向量，考虑这个方程组。
           \begin{equation*}
             \begin{linsys}{2}
               (3-x)\cdot b_1  &+  &2\cdot b_2  &=  &0  \\
               -1\cdot b_1     &-  &x\cdot b_2  &=  &0
             \end{linsys}
           \end{equation*}
           取 $\lambda_1=2$ 我们得到
           \begin{equation*}
             \begin{linsys}{2}
                1\cdot b_1  &+  &2\cdot b_2  &=  &0  \\
               -1\cdot b_1  &-  &2\cdot b_2  &=  &0
             \end{linsys}
           \end{equation*}
           由此得到这个特征子空间和特征向量。
           \begin{equation*}
             \set{\colvec{-2b_2 \\ b_2}
                   \suchthat b_2\in\C}
             \qquad
             \colvec[r]{-2 \\ 1}
           \end{equation*}
           取 $\lambda_2=1$，方程组为
           \begin{equation*}
             \begin{linsys}{2}
                2\cdot b_1  &+  &2\cdot b_2  &=  &0  \\
               -1\cdot b_1  &-  &1\cdot b_2  &=  &0
             \end{linsys}
           \end{equation*}
           由此得到这个。
           \begin{equation*}
             \set{\colvec{-b_2 \\ b_2}
                   \suchthat b_2\in\C}
             \qquad
             \colvec[r]{-1 \\ 1}
           \end{equation*}
       \end{exparts}  
     \end{answer}
  \item
     求这个矩阵的特征方程以及特征值和相应的特征向量。
     \textit{提示。}
       特征值是复数。
     \begin{equation*}
       \begin{mat}[r]
         -2  &-1 \\
          5  &2
      \end{mat}  
    \end{equation*}
    \begin{answer}
         特征方程
           \begin{equation*}
             0=
             \begin{vmatrix}
               -2-x  &-1  \\
               5     &2-x               
             \end{vmatrix}
             =x^2+1
           \end{equation*}
           有复根 $\lambda_1=i$ 和 $\lambda_2=-i$。
           对于方程组
           \begin{equation*}
             \begin{linsys}{2}
               (-2-x)\cdot b_1  &-  &1\cdot b_2      &=  &0  \\
               5\cdot b_1       &   &(2-x)\cdot b_2  &=  &0
             \end{linsys}
           \end{equation*}
           取 $\lambda_1=i$，用高斯消元法得到如下化简。
           \begin{equation*}
             \begin{linsys}{2}
                (-2-i)\cdot b_1  &-  &1\cdot b_2      &=  &0  \\
                5\cdot b_1       &-  &(2-i)\cdot b_2  &=  &0
             \end{linsys}
             \grstep{(-5/(-2-i))\rho_1+\rho_2}
             \begin{linsys}{2}
                (-2-i)\cdot b_1  &-  &1\cdot b_2      &=  &0  \\
                                 &   &0               &=  &0
             \end{linsys}
           \end{equation*}
           （对于右下角的计算，通分后取公分母
           \begin{equation*}
             \frac{5}{-2-i}-(2-i)
             =
             \frac{5}{-2-i}-\frac{-2-i}{-2-i}\cdot (2-i)
             =
             \frac{5-(-5)}{-2-i}
           \end{equation*}
           即可看出它给出一个 $0=0$ 的方程。）
           由此得到特征子空间和特征向量。
           \begin{equation*}
             \set{\colvec{(1/(-2-i))b_2 \\ b_2}
                   \suchthat b_2\in\C}
             \qquad
             \colvec{1/(-2-i) \\ 1}
           \end{equation*}
           取 $\lambda_2=-i$，方程组
           \begin{equation*}
             \begin{linsys}{2}
                (-2+i)\cdot b_1  &-  &1\cdot b_2      &=  &0  \\
                5\cdot b_1       &-  &(2+i)\cdot b_2  &=  &0
             \end{linsys}
             \grstep{(-5/(-2+i))\rho_1+\rho_2}
             \begin{linsys}{2}
                (-2+i)\cdot b_1  &-  &1\cdot b_2      &=  &0  \\
                                 &   &0               &=  &0
             \end{linsys}
           \end{equation*}
           导出这个。
           \begin{equation*}
             \set{\colvec{(1/(-2+i))b_2 \\ b_2}
                   \suchthat b_2\in\C}
             \qquad
             \colvec{1/(-2+i) \\ 1}
           \end{equation*}
    \end{answer}
  \item  
    求这个矩阵的特征多项式、特征值以及相应的特征向量。
    \begin{equation*}
      \begin{mat}[r]
        1  &1  &1  \\
        0  &0  &1  \\
        0  &0  &1
      \end{mat}
    \end{equation*}
    \begin{answer}
      特征方程是
      \begin{equation*}
        0=
        \begin{vmatrix}
          1-x  &1   &1   \\
          0    &-x  &1   \\
          0    &0   &1-x
        \end{vmatrix}
        =(1-x)^2(-x)
      \end{equation*}
      因此特征值为 $\lambda_1=1$（它是方程的一个重根）和 $\lambda_2=0$。
      对于其余部分，考虑这个方程组。
      \begin{equation*}
        \begin{linsys}{3}
          (1-x)\cdot b_1  &+  &b_2         &+  &b_3            &=  &0  \\
                          &   &-x\cdot b_2 &+  &b_3            &=  &0  \\
                          &   &            &   &(1-x)\cdot b_3 &= &0  
        \end{linsys}
      \end{equation*}
      当 $x=\lambda_1=1$ 时，解集是这个特征子空间。
      \begin{equation*}
        \set{\colvec{b_1 \\ 0 \\ 0}\suchthat b_1\in\C}
      \end{equation*}
      当 $x=\lambda_2=0$ 时，解集是这个特征子空间。
      \begin{equation*}
        \set{\colvec{-b_2 \\ b_2 \\ 0}\suchthat b_2\in\C}
      \end{equation*}
      所以这些就是与 $\lambda_1=1$ 和 
      $\lambda_2=0$ 相对应的特征向量。
      \begin{equation*}
        \colvec[r]{1 \\ 0 \\ 0}
        \qquad
        \colvec[r]{-1 \\ 1 \\ 0}  
      \end{equation*}
    \end{answer}
  \recommended \item
    对每个矩阵，求出特征方程以及特征值和相应的特征向量。
    \begin{exparts*}
      \partsitem \( \begin{mat}[r]
              3  &-2 &0  \\
             -2  &3  &0  \\
              0  &0  &5
            \end{mat}  \)
      \partsitem \( \begin{mat}[r]
              0  &1   &0  \\
              0  &0   &1  \\
              4  &-17 &8
            \end{mat}  \)
    \end{exparts*}
    \begin{answer}
      \begin{exparts}
         \partsitem 特征方程是
           \begin{equation*}
             0=
             \begin{vmatrix}
              3-x  &-2   &0  \\
             -2    &3-x  &0  \\
              0    &0    &5-x
             \end{vmatrix}
             =x^3-11x^2+35x-25=(x-1)(x-5)^2
           \end{equation*}
           因此特征值为 $\lambda_1=1$，还有重复的特征值
           $\lambda_2=5$。
           为求特征向量，考虑这个方程组。
           \begin{equation*}
             \begin{linsys}{3}
               (3-x)\cdot b_1  &-  &2\cdot b_2      &   &   &=  &0  \\
               -2\cdot b_1     &+  &(3-x)\cdot b_2  &   &   &=  &0  \\
                               &   &                &   &(5-x)\cdot b_3 &= &0 
             \end{linsys}
           \end{equation*}
           取 $\lambda_1=1$ 我们得到
           \begin{equation*}
             \begin{linsys}{3}
                2\cdot b_1     &-  &2\cdot b_2   &   &   &=  &0  \\
               -2\cdot b_1     &+  &2\cdot b_2   &   &   &=  &0  \\
                               &   &             &   &4\cdot b_3 &= &0 
             \end{linsys}
           \end{equation*}
           由此得到这个特征子空间和特征向量。
           \begin{equation*}
             \set{\colvec{b_2 \\ b_2 \\ 0}
                   \suchthat b_2\in\C}
             \qquad
             \colvec[r]{1 \\ 1 \\ 0}
           \end{equation*}
           取 $\lambda_2=5$，方程组为
           \begin{equation*}
             \begin{linsys}{3}
                -2\cdot b_1  &-  &2\cdot b_2   &   &   &=  &0  \\
               -2\cdot b_1   &-  &2\cdot b_2   &   &   &=  &0  \\
                             &   &             &   &0\cdot b_3 &= &0 
             \end{linsys}
           \end{equation*}
           由此得到这个。
           \begin{equation*}
             \set{\colvec{-b_2 \\ b_2 \\ 0}+\colvec{0 \\ 0 \\ b_3}
                   \suchthat b_2,b_3\in\C}
             \qquad
             \colvec[r]{-1 \\ 1 \\ 0},\,\colvec[r]{0 \\ 0 \\ 1}
           \end{equation*}
         \partsitem 特征方程是
           \begin{equation*}
             0=
             \begin{vmatrix}
              -x   &1    &0  \\
              0    &-x   &1  \\
              4    &-17  &8-x
             \end{vmatrix}
             =-x^3+8x^2-17x+4=-1\cdot(x-4)(x^2-4x+1)
           \end{equation*}
           特征值为 $\lambda_1=4$ 以及（用求根公式）
           $\lambda_2=2+\sqrt{3}$ 和 
           $\lambda_3=2-\sqrt{3}$。
           为求特征向量，考虑这个方程组。
           \begin{equation*}
             \begin{linsys}{3}
               -x\cdot b_1  &+  &b_2          &   &               &=  &0  \\
                            &   &-x\cdot b_2  &+  &b_3            &=  &0  \\
               4\cdot b_1   &-  &17\cdot b_2  &+  &(8-x)\cdot b_3 &= &0 
             \end{linsys}
           \end{equation*}
           代入 $x=\lambda_1=4$ 得到方程组
           \begin{align*}
             \begin{linsys}{3}
               -4\cdot b_1  &+  &b_2          &   &           &= &0  \\
                            &   &-4\cdot b_2  &+  &b_3        &= &0  \\
               4\cdot b_1   &-  &17\cdot b_2  &+  &4\cdot b_3 &= &0 
             \end{linsys}                                              
             &\grstep{\rho_1+\rho_3}
             \begin{linsys}{3}
               -4\cdot b_1  &+  &b_2          &   &           &= &0  \\
                            &   &-4\cdot b_2  &+  &b_3        &= &0  \\
                            &   &-16\cdot b_2  &+ &4\cdot b_3 &= &0 
             \end{linsys}                                                \\
             &\grstep{-4\rho_2+\rho_3}
             \begin{linsys}{3}
               -4\cdot b_1  &+  &b_2          &   &           &= &0  \\
                            &   &-4\cdot b_2  &+  &b_3        &= &0  \\
                            &   &              &  &0          &= &0 
             \end{linsys}
           \end{align*}
           由此得到这个特征子空间和特征向量。
           \begin{equation*}
             V_4=\set{\colvec{(1/16)\cdot b_3 \\ (1/4)\cdot b_3 \\ b_3}
                   \suchthat b_2\in\C}
             \qquad
             \colvec[r]{1 \\ 4 \\ 16}
           \end{equation*}

           代入 $x=\lambda_2=2+\sqrt{3}$ 得到方程组
           \begin{multline*}
             \begin{linsys}{3}
               (-2-\sqrt{3})\cdot b_1  
                      &+  &b_2          
                          &   &           &= &0  \\
                      &   &(-2-\sqrt{3})\cdot b_2  
                          &+  &b_3        &= &0  \\
               4\cdot b_1   
                      &-  &17\cdot b_2  
                          &+  &(6-\sqrt{3})\cdot b_3 &= &0 
             \end{linsys}                                   \\
             \grstep{(-4/(-2-\sqrt{3}))\rho_1+\rho_3}
             \begin{linsys}{3}
               (-2-\sqrt{3})\cdot b_1  
                      &+  &b_2          
                          &   &           &= &0  \\
                      &   &(-2-\sqrt{3})\cdot b_2  
                          &+  &b_3        &= &0  \\
                      &+  &(-9-4\sqrt{3})\cdot b_2  
                          &+  &(6-\sqrt{3})\cdot b_3 &= &0 
             \end{linsys}
           \end{multline*}
           （第三个方程的中间系数等于数 $(-4/(-2-\sqrt{3}))-17$；
           先取 $-2-\sqrt{3}$ 为公分母，
           再将分数的分子分母同乘 $-2+\sqrt{3}$ 以使分母有理化）
           \begin{equation*}
             \grstep{((9+4\sqrt{3})/(-2-\sqrt{3}))\rho_2+\rho_3}
             \begin{linsys}{3}
               (-2-\sqrt{3})\cdot b_1  
                      &+  &b_2          
                          &   &           &= &0  \\
                      &   &(-2-\sqrt{3})\cdot b_2  
                          &+  &b_3        &= &0  \\
                      &   &                         
                          &   &0                     &= &0 
             \end{linsys}
           \end{equation*}
           由此得到这个特征子空间和特征向量。
           \begin{equation*}
             V_{2+\sqrt{3}}
             =\set{\colvec{(1/(2+\sqrt{3})^2)\cdot b_3  \\ 
                           (1/(2+\sqrt{3}))\cdot b_3    \\ 
                           b_3}
                    \suchthat b_3\in\C}
             \qquad
             \colvec{(1/(2+\sqrt{3})^2)  \\ 
                           (1/(2+\sqrt{3}))  \\ 
                           1}
           \end{equation*}

           最后，代入 $x=\lambda_3=2-\sqrt{3}$ 得到方程组
           \begin{multline*}
             \begin{linsys}{3}
               (-2+\sqrt{3})\cdot b_1  
                      &+  &b_2          
                          &   &           &= &0  \\
                      &   &(-2+\sqrt{3})\cdot b_2  
                          &+  &b_3        &= &0  \\
               4\cdot b_1   
                      &-  &17\cdot b_2  
                          &+  &(6+\sqrt{3})\cdot b_3 &= &0 
             \end{linsys}                                       \\
             \begin{aligned}
               &\grstep{(-4/(-2+\sqrt{3}))\rho_1+\rho_3}
                \begin{linsys}{3}
                  (-2+\sqrt{3})\cdot b_1  
                         &+  &b_2          
                             &   &           &= &0  \\
                         &   &(-2+\sqrt{3})\cdot b_2  
                             &+  &b_3        &= &0  \\
                         &   &(-9+4\sqrt{3})\cdot b_2  
                             &+  &(6+\sqrt{3})\cdot b_3 &= &0 
                \end{linsys}                                       \\
                &\grstep{((9-4\sqrt{3})/(-2+\sqrt{3}))\rho_2+\rho_3}
                \begin{linsys}{3}
                  (-2+\sqrt{3})\cdot b_1  
                         &+  &b_2          
                             &   &           &= &0  \\
                         &   &(-2+\sqrt{3})\cdot b_2  
                             &+  &b_3        &= &0  \\
                         &   &                         
                             &   &0                     &= &0 
                \end{linsys}
              \end{aligned}
           \end{multline*}
           由此得到这个特征子空间和特征向量。
           \begin{equation*}
             V_{2-\sqrt{3}}
             =\set{\colvec{(1/(-2+\sqrt{3})^2)\cdot b_3  \\ 
                           (1/(2-\sqrt{3}))\cdot b_3    \\ 
                           b_3}
                    \suchthat b_3\in\C}
             \qquad
             \colvec{(1/(-2+\sqrt{3})^2)  \\ 
                           (1/(2-\sqrt{3}))  \\ 
                           1}
           \end{equation*}
       \end{exparts}
     \end{answer}
  \item 
   对每个矩阵，求出特征多项式以及特征值和相应的特征子空间。
   同时求出代数重数与几何重数。
    \begin{exparts*}
      \partsitem 
        $\begin{mat}
           13 &-4 \\
           -4 &7
         \end{mat}$
      \partsitem
         $\begin{mat}
           1 &3 &-3 \\
          -3 &7 &-3 \\
          -6 &6 &-2
         \end{mat}$
     \partsitem
         $\begin{mat}
           2 &3 &-3 \\
           0 &2 &-3 \\
           0 &0 &1
         \end{mat}$
    \end{exparts*}
    \begin{answer}
      \begin{exparts}
        \item 特征多项式可分解为 
          $x^2-20x+75=(x-5)(x-15)$，
          因此特征值为 $\lambda_1=5$ 和~$\lambda_2=15$。
          下面是相应的特征子空间。
          \begin{equation*}
            V_5=\set{k\colvec{1 \\ 2}\suchthat k\in\C}
            \quad
            V_{15}=\set{k\colvec{-2 \\ 1}\suchthat k\in\C}
          \end{equation*}
          对每个特征值，代数重数与几何重数都是~$1$。
       \item 特征多项式 $x^3 - 6x^2 + 32$ 可分解为
          $(x+2)(x-4)^2$。
          特征向量为 $\lambda_1=-2$ 和 $\lambda_2=4$。
          下面是相应的特征子空间。
          \begin{equation*}
            V_{-2}=\set{k\colvec{1 \\ 1 \\ 2}\suchthat k\in\C}
            \quad
            V_4=\set{k_1\colvec{1 \\ 1 \\ 0}
                      +k_2\colvec{1 \\ 0 \\ -1} \suchthat k_1,k_2\in\C}
          \end{equation*}
          对于 $\lambda_1=-2$，代数重数与几何重数
          都是~$1$。 
          对于 $\lambda_2=4$，代数重数与几何重数
          都是~$2$。
       \item 特征多项式 $x^3 - 5x^2 + 8x - 4$ 可分解为
          $(x-1)(x-2)^2$。
          特征值为 $\lambda_1=1$ 和 $\lambda_2=2$。
          下面是相应的特征子空间。
          \begin{equation*}
            V_1=\set{k\colvec{-6 \\ 3 \\ 1}\suchthat k\in \C}
            \quad
            V_2=\set{k\colvec{1 \\ 0 \\ 0}\suchthat k\in\C}
          \end{equation*}
          对于 $\lambda_1=1$，代数重数与几何重数
          都是~$1$。
          对于 $\lambda_2=2$，代数重数是~$2$，但
          几何重数是~$1$。
      \end{exparts}
    \end{answer}
   \recommended \item
     设 \( \map{t}{\polyspace_2}{\polyspace_2} \) 是这个线性映射。
     \begin{equation*}
      a_0+a_1x+a_2x^2\mapsto
      (5a_0+6a_1+2a_2)-(a_1+8a_2)x+(a_0-2a_2)x^2
     \end{equation*}
    求出它的特征值和相应的特征向量。
    \begin{answer}
      关于自然基 $B=\sequence{1,x,x^2}$ 
      矩阵表示如下。
      \begin{equation*}
        \rep{t}{B,B}
        =
        \begin{mat}[r]
          5  &6  &2  \\
          0  &-1 &-8 \\
          1  &0  &-2 
        \end{mat}
      \end{equation*}
      于是特征方程
      \begin{equation*}
        0
        =
        \begin{mat}
          5-x  &6    &2  \\
          0    &-1-x &-8 \\
          1    &0    &-2-x 
        \end{mat}
        =(5-x)(-1-x)(-2-x)-48-2\cdot(-1-x)
      \end{equation*}
      为 $0=-x^3+2x^2+15x-36=-1\cdot (x+4)(x-3)^2$。
      为求相应的特征向量，考虑这个方程组。
      \begin{equation*}
        \begin{linsys}{3}
          (5-x)\cdot b_1 &+ &6b_2      &+ &2b_3      &= &0 \\
                         &  &(-1-x)\cdot b_2 &- &8b_3      &= &0 \\
          b_1            &  &                &+ &(-2-x)\cdot b_3 &= &0 
        \end{linsys}
      \end{equation*}
      将 $\lambda_1=-4$ 代入 $x$ 得到
      \begin{equation*}
